\begin{statementbox}
\paragraph{极值点：局部比较，不是全局比较}
设 $x_0$ 是定义域的内点。若存在 $\delta>0$，使 $0<|x-x_0|<\delta$ 时均有
$f(x)<f(x_0)$，则 $x_0$ 是严格极大值点，$f(x_0)$ 是严格极大值；将 $<$ 换成 $>$ 就得到严格极小值。
若允许等号，则称非严格的局部极值。下文的变号判定给出严格极值。

\textbf{必要条件（费马定理）：}若 $f$ 在 $x_0$ 处可导且取得局部极值，则 $f'(x_0)=0$；反过来不成立。

\textbf{一阶导数变号判定：}若 $f$ 在 $x_0$ 连续、在其两侧某去心邻域可导，则
\begin{itemize}
\item 左侧 $f'>0$、右侧 $f'<0$，则 $x_0$ 为严格极大值点；
\item 左侧 $f'<0$、右侧 $f'>0$，则 $x_0$ 为严格极小值点；
\item 两侧都 $f'>0$，或两侧都 $f'<0$，则 $x_0$ 不是极值点。
\end{itemize}
变号判定是充分判据；未出现上述稳定符号情形时，须另行判断，不能直接说没有极值。
\end{statementbox}
